\section{Instance-Level Semantics for Constraint-Aware Counterfactual Query}

% \input{section/4_instance/table/modality}
\label{sec:instance}
Given a query $q$, the system retrieves instances in $f_\omega$ that satisfy the instance-level constraints and remain competitive under the cost function $c_\omega$. This section analyzes how P1 (traceability), P2 (cost faithfulness), and P3 (efficiency) manifest across different data modalities, highlighting where admissibility can be verified and how constraint evaluation shapes system-level trade-offs in counterfactual retrieval. The set-level collective validity P4 is discussed separately in Section~\ref{sec:set}.


\paragraph{Tabular and Relational Data.}
Among all modalities, tabular data exhibits the most explicit instance-level structure. Because features are represented column-wise, modifications have clear coordinate semantics. The interpretable edit $\Delta(x_q,x)$ can be constructed directly through column-wise comparison (P1), the transformation cost $\omega(x_q,x)$ can aggregate per-column modification effort in a semantically meaningful way (P2), and immutable attributes can be encoded directly within $\omega$, where infeasible changes are assigned infinite cost. In this sense, instance-level semantics are structurally transparent in tabular domains and align naturally with P1 and P2.

Despite this structural transparency, several design challenges remain. (i) \emph{Heterogeneous edit costs}~\cite{karabulut2019weighted}. Different columns carry distinct semantic scales and actionability, and continuous and categorical attributes coexist. A naive geometric metric may therefore violate cost alignment, making the design of $\omega$ a modeling problem rather than a purely geometric one. (ii) \emph{Constraint-induced partitioning}~\cite{galhotra2022hyper}. Immutable attributes and structural constraints partition the admissible space induced by $\omega$ into logically defined subspaces. Index construction must respect these partitions; otherwise, geometrically close candidates may lie outside the admissible region, leading to inefficient validation overhead. (iii) \emph{Ranking–admissibility decoupling}~\cite{bartal2011fast}. Even after partitioning, geometric distance provides only an ordering heuristic. The transformation cost $\omega(x_q,x)$ induces a ranking over candidates, while admissibility is implicitly determined by whether $\omega(x_q,x) < \infty$. Distance alone cannot guarantee feasibility. Without compiling admissibility constraints into the index, the system may repeatedly explore top-ranked but infeasible candidates, shifting the computational burden from metric navigation to constraint evaluation. Efficient retrieval therefore requires integrating column-level semantics into the index structure to ensure scalable access to viable candidates (P3).

\paragraph{Sequence and Temporal Data.}
In sequential domains, the query instance $x_q$ is represented as a trajectory $(x_{q,1},\dots,x_{q,T})$. Measuring transformation cost therefore requires temporal alignment, and the transformation cost $\omega(x_q,x)$ typically relies on dynamic time warping (DTW)~\cite{senin2008dynamic} or learned sequence embeddings~\cite{nalmpantis2019signal2vec}. DTW captures morphological similarity by optimizing over time mappings, providing a geometric notion of proximity. However, temporal structure introduces constraints that are defined by the underlying dynamics, not by the aligned observations. For example, in chess, two sequences of moves may be close under an alignment of move tokens, yet the aligned edits can imply an illegal continuation because each move depends on the current board state. Accordingly, admissibility under $\omega$ depends on the execution induced by the sequence. As a result, the transformation cost and admissibility induced by $\omega$ may reside in different semantic layers: similarity is computed in observation space, while feasibility is determined in execution space.

This structural separation appears in three dimensions. (i) \emph{Execution-aware edit representation}~\cite{jing2024towards}. In sequential domains, edits cannot be reduced to independent coordinate substitutions. They must correspond to executable transformations under system dynamics. Local observation-level deviations may lack a coherent execution interpretation, especially when early edits alter downstream states. Constructing $\Delta(x_q,x)$ therefore requires aligning deviations with valid execution paths, elevating representation from coordinate differences to trajectory-level transformations. Beyond DTW, optimal transport~\cite{charpentier2023optimal} models global correspondences under structural constraints, enabling semantically coherent edits (P1). (ii) \emph{Temporal sensitivity}~\cite{rajput2024evaluating}. Small time shifts or rescaling can substantially change $\omega(x_q,x)$ even when behaviors are similar. Signal-space metrics such as DTW may misestimate intervention cost, complicating faithful modeling (P2). (iii) \emph{Path dependency and state propagation}~\cite{page2006path}. In stateful systems, early edits propagate and affect downstream states. Admissibility under $\omega$ therefore requires validation under system dynamics rather than step-wise substitution, shifting computation toward execution-level verification and impacting scalable retrieval (P3). Efficient designs may combine alignment lower bounds with checkpoints or constraint automata to reduce repeated validation.

\paragraph{Graph and Network Data.}

In graph and network domains, the query instance $x_q$ and a candidate $x$ are structured objects composed of nodes and edges. Let $\Delta(x_q,x)$ denote the edit set transforming $x_q$ into $x$. Measuring transformation cost is computationally challenging, as exact graph edit distance~\cite{gao2010survey} is expensive and unstable under small perturbations; therefore, the transformation cost $\omega(x_q,x)$ is typically approximated in an embedding space~\cite{kipf2016semi}, such as $\|\phi(x_q)-\phi(x)\|_2$. Embeddings provide scalable geometric proximity, yet structural legality is determined by the constraints encoded in $\omega$, which require $\Delta(x_q,x)$ to satisfy application-specific constraints and structural invariants such as connectivity, typing, capacity, or separation rules. Similar to the temporal case, similarity and admissibility operate over distinct semantic layers.

Graph-structured data introduces three sources of complexity. (i) \emph{Multi-level semantics}~\cite{liang2021mile}. Edits may occur at the level of nodes, edges, subgraphs, or entire graphs, each with different semantic implications. A local modification can induce a substantial structural shift depending on granularity, complicating traceable edit representation (P1). (ii) \emph{Embedding–constraint misalignment}~\cite{armgaan2024graphtrail}. Embedding space captures predictive similarity but may align weakly with structural invariants, so candidates ranked highly under $\omega(x_q,x)$ may violate structural constraints encoded in $\omega$, undermining cost faithfulness (P2). (iii) \emph{Invariant verification cost}~\cite{xu2018powerful}. Many graph constraints are global; verifying them may require whole-graph computation and become the dominant runtime cost as graph size grows, impacting scalable retrieval (P3). Practical systems can maintain incremental summaries of invariants, restrict admissible edit spaces, or precompute locality-aware metadata to reduce reconstruction overhead for $\Delta(x_q,x)$ while preserving structural correctness.

\paragraph{Text.}
In text domains, the query instance $x_q$ is a token sequence that carries multiple layers of structure, including syntax, discourse, and factual semantics. Let $\Delta(x_q,x)$ denote the token- or phrase-level edit transformation from $x_q$ to a candidate text $x$. Relying solely on surface metrics such as edit distance or BLEU~\cite{papineni2002bleu} to measure transformation cost fails to capture semantic equivalence and contextual meaning, and therefore cannot faithfully reflect true semantic change. In practice, the transformation cost $\omega(x_q,x)$ is typically approximated in a pretrained representation space~\cite{devlin2019bert} to enable scalable semantic similarity estimation. However, textual admissibility is determined by constraints encoded in $\omega$, which require $\Delta(x_q,x)$ to preserve semantic coherence and satisfy conditions such as entailment, factual consistency, or policy compliance.

This separation between representation space and semantic constraints introduces three sources of complexity. (i) \emph{Granularity mismatch}~\cite{yin2015multigrancnn}. Edits occur at the token or phrase level, yet their effects propagate to sentence-level intent, discourse relations, and factual stance. A local change may alter presuppositions or coreference, making the mapping from $\Delta(x_q,x)$ to semantic effect non-linear and context-dependent, and complicating traceable edit representation (P1). (ii) \emph{Similarity–validity misalignment}~\cite{wang2021robustness}. Representation proximity captures distributional similarity but aligns weakly with entailment, factual correctness, or policy compliance. Candidates ranked highly under $\omega(x_q,x)$ may still violate semantic constraints encoded in $\omega$, showing that geometric similarity does not ensure admissibility and challenging cost alignment (P2). (iii) \emph{Context-dependent verification}~\cite{gu2024survey}. Because semantic effects extend beyond local edits, admissibility cannot be checked at the token level alone. Feasibility under $\omega$ often requires full-sequence reasoning or external grounding, turning constraint evaluation into a global inference problem that can dominate runtime and affect scalable retrieval (P3). Practical systems therefore adopt multi-stage designs that combine representation-based retrieval with layered semantic validation, using cached annotations or incremental evaluation to reduce repeated verification.

\paragraph{Image.}
Among all modalities, image data poses the greatest challenge for interpretable editing. An image is a high-dimensional pixel array mapped by a feature extractor $\phi$ into an embedding space for scalable retrieval. Yet pixel-level changes lack intrinsic semantic meaning: small perturbations may cause large embedding shifts without real object-level edits, while subtle visual differences can alter task-critical attributes~\cite{narodytska2017simple}. Feasibility is defined at the level of objects, regions, or scene structure rather than pixels. Thus, admissibility under $\omega$ must enforce object-level transformations that satisfy domain constraints such as anatomical validity or physical plausibility~\cite{nam2026causal}, separating similarity from admissibility across distinct semantic layers.

Despite the scalability of embedding-based image retrieval, several challenges remain. (i) \emph{Semantic opacity of pixel edits}~\cite{assran2023self}. Pixel differences do not directly correspond to object- or attribute-level changes, making reconstruction of a traceable edit $\Delta(x_q,x)$ harder than in structured domains (P1). (ii) \emph{Perceptual–task inconsistency}~\cite{zhang2021multi}. Representation proximity reflects perceptual similarity but may overlook task-critical regions; visually similar images can differ in clinically or functionally decisive ways, challenging cost faithfulness (P2). (iii) \emph{Region-aware admissibility checking}~\cite{gagliardi2025region}. Feasibility under $\omega$ often depends on localization, segmentation, or scene reasoning, requiring auxiliary vision modules that can dominate runtime and limit scalable querying (P3). Recent efforts introduce intermediate semantic layers, such as concept bottleneck models~\cite{de2025causally,koh2020concept,dominici2025counterfactual}, which transform images into interpretable concept vectors prior to prediction. Operating in concept space improves cost faithfulness and enables constraints on meaningful attributes, but it demands extensive concept annotation and can still be difficult to map changes back to the original input images.


\paragraph{Extended to Multi-Modality.}
When instances span multiple modalities, counterfactual querying becomes a coordinated validation problem in which feasibility is defined conjunctively across modality-specific predicates together with cross-modal consistency constraints~\cite{chen2019cross}. A candidate must satisfy each structural or semantic rule imposed by its modality while remaining mutually coherent across representations, so images, text, and structured attributes must be individually valid and jointly aligned. As modalities accumulate, admissibility complexity grows combinatorially, and similarity search, whether performed in joint or modality-specific embedding spaces, functions primarily as a proposal mechanism. Correctness is ultimately determined through structured cross-modal certification that validates each modality-specific constraint and enforces agreement among them. The dominant computational burden therefore shifts from geometric navigation to coordinated validation, making scalable multi-modal counterfactual systems fundamentally verification-driven rather than similarity-driven.
